\section{Введение}

\subsection{Контекст и актуальность}

В предыдущей научно-исследовательской работе была воспроизведена фрактальная
граница обучаемости нейронной сети в двумерном пространстве скоростей обучения
слоёв $(\eta_0,\eta_1)$, измерена box-counting размерность границы для активаций
tanh, ReLU и линейной сети, а также показана связь наблюдаемой картины с режимами
динамики градиентного спуска: монотонной сходимостью, catapult-режимом,
периодичностью, хаосом и расходимостью.

Ограничение того исследования очевидно: настоящая практика обучения нейронных
сетей почти никогда не сводится к чистому градиентному спуску с двумя шагами
обучения. Реальный оптимизатор содержит инерцию (momentum), регуляризацию
затуханием весов (weight decay), стохастичность мини-батчей, а на практике чаще
всего применяются адаптивные методы --- Adam и RMSprop. Поэтому естественный и
содержательный вопрос состоит в том, является ли фрактальность границы
обучаемости свойством именно «голого» градиентного спуска, или же она сохраняется
и в более реалистичном многомерном пространстве гиперпараметров.

Практическая значимость вопроса связана с автоматическим подбором
гиперпараметров. Если граница между успешным и неуспешным обучением остаётся
фрактальной при добавлении инерции и регуляризации, то любые методы поиска,
опирающиеся на гладкость или липшицевость целевой метрики по гиперпараметрам,
принципиально ограничены: на любом масштабе рядом с удачной конфигурацией
находится неудачная. Если же адаптивные методы сглаживают границу, это даёт
количественный аргумент в пользу их робастности, помимо привычных аргументов
о скорости сходимости.

\subsection{Цель и задачи работы}

\textbf{Цель работы:} исследовать, сохраняется ли фрактальная структура границы
обучаемости нейронной сети при переходе от двумерного пространства скоростей
обучения к пространству, включающему инерцию, затухание весов, размер батча и
параметры адаптивных оптимизаторов, и количественно сравнить сложность границы
для разных методов оптимизации.

\textbf{Задачи:}
\begin{enumerate}[leftmargin=*,topsep=2pt,itemsep=2pt]
  \item Кратко изложить результаты предыдущей НИР и обзор литературы по
        фрактальной границе обучаемости, эффекту грани устойчивости и роли
        крупных шагов обучения.
  \item Получить теоретический мини-результат: для квадратичной функции
        $L(w)=\tfrac12\lambda w^2$ записать градиентный спуск с инерцией в виде
        двумерного линейного отображения, найти условия локальной устойчивости
        через собственные значения матрицы перехода и сравнить область
        устойчивости в плоскости $(\eta,\beta)$ с обычным градиентным спуском.
  \item Воспроизвести базовый эксперимент предыдущей НИР и построить новые
        двумерные сечения пространства гиперпараметров: $(\eta,\beta)$,
        $(\eta,\lambda_{wd})$, $(\eta,B)$, с цветовой картой режимов и оценкой
        box-counting размерности границы.
  \item Сравнить SGD, SGD с инерцией, Adam и RMSprop; проверить гипотезу о том,
        что адаптивные методы сглаживают фрактальную структуру границы за счёт
        покоординатной нормализации градиентов.
  \item Сравнить границы обучаемости на случайных и на структурированных данных
        (набор \texttt{sklearn digits}).
  \item Построить трёхмерный эксперимент $(\eta,\beta,\lambda_{wd})$ в виде
        набора двумерных сечений и обсудить, можно ли по сечениям судить о
        сложности границы в полном трёхмерном пространстве.
  \item Обеспечить воспроизводимость: фиксация seed, типа точности, разрешения
        сетки и критериев классификации режимов; оформление отчёта в \LaTeX.
\end{enumerate}

\subsection{Терминология}

Далее используются следующие соответствия терминов:
\emph{trainability boundary} --- граница обучаемости;
\emph{box-counting dimension} --- размерность Минковского, или box-counting
размерность; \emph{learning rate} --- шаг обучения; \emph{momentum} --- инерция;
\emph{weight decay} --- затухание весов, или $L_2$-регуляризация;
\emph{edge of stability} --- грань устойчивости; \emph{catapult effect} ---
catapult-эффект, или режим выброса.

\subsection{Структура работы}

В разделе~\ref{sec:problem} даётся строгая постановка задачи: модель, семейство
оптимизаторов, критерии классификации режимов обучения и способ измерения
размерности. Раздел~\ref{sec:review} содержит обзор литературы.
Раздел~\ref{sec:theory} содержит теоретический результат об области устойчивости
метода тяжёлого шарика. Раздел~\ref{sec:methods} описывает программную
реализацию и вычислительный протокол. Раздел~\ref{sec:results} представляет
результаты экспериментов, раздел~\ref{sec:discussion} --- их обсуждение,
связь с теорией и ограничения. Раздел~\ref{sec:conclusion} подводит итоги.
